% 三视图横向并排点修：投影边线/尺寸/虚实均沿用既有稿，仅平移scope使三幅同行，图名基线一致
\begin{tikzpicture}[line width=.6pt,font=\normalsize,x=1cm,y=1cm]
% 正视图：矩形(0,0)-(3,3)表示6cm，虚线x=1(2cm处)，实线x=2(4cm处)
\begin{scope}[shift={(0,0)}]
\draw (0,0) rectangle (3,3);
\draw[dashed] (1,0) -- (1,3);
\draw (2,0) -- (2,3);
\draw[<->] (0,-0.35) -- (1,-0.35);
\node[anchor=north,font=\normalsize] at (0.5,-0.4) {2};
\draw[<->] (1,-0.35) -- (2,-0.35);
\node[anchor=north,font=\normalsize] at (1.5,-0.4) {2};
\draw[<->] (2,-0.35) -- (3,-0.35);
\node[anchor=north,font=\normalsize] at (2.5,-0.4) {2};
\draw[<->] (3.35,0) -- (3.35,3);
\node[anchor=west,font=\normalsize] at (3.4,1.5) {6};
\node[below] at (1.5,-0.95) {正视图};
\end{scope}
% 侧视图：矩形(0,0)-(3,3)表示深6cm，实线x=1.5(3cm处)
\begin{scope}[shift={(4.6,0)}]
\draw (0,0) rectangle (3,3);
\draw (1.5,0) -- (1.5,3);
\draw[<->] (0,-0.35) -- (1.5,-0.35);
\node[anchor=north,font=\normalsize] at (0.75,-0.4) {3};
\draw[<->] (1.5,-0.35) -- (3,-0.35);
\node[anchor=north,font=\normalsize] at (2.25,-0.4) {3};
\node[below] at (1.5,-0.95) {侧视图};
\end{scope}
% 俯视图：五边形 (0,1.5)-(0,3)-(3,3)-(3,0)-(1.5,0)，外包3x3=正视宽x侧视宽
\begin{scope}[shift={(9.2,0)}]
\draw (0,1.5) -- (0,3) -- (3,3) -- (3,0) -- (1.5,0) -- cycle;
\node[below] at (1.5,-0.95) {俯视图};
\end{scope}
\end{tikzpicture}
